\uuid{aUn2}
\chapitre{Statistique}
\niveau{L2}
\module{Analyse de données}
\sousChapitre{Tests d'hypothèses, intervalle de confiance}
\titre{Contrôle qualité}
\theme{statistiques, tests d'hypothèses}
\auteur{}
\datecreate{2022-09-28}
\organisation{AMSCC}
\difficulte{2}
\contenu{
\texte{On contrôle 50 épaisseurs indépendantes de même distribution, supposée de variance finie. Les valeurs observées, en mm, sont :
\begin{center}\begin{tabular}{lrrrr}\hline Épaisseur & 4,8 & 4,9 & 5,0 & 5,1 \\
Effectif & 5 & 15 & 20 & 10 \\ \hline\end{tabular}\end{center}
On utilise ici l'approximation normale de la moyenne pour un échantillon de grande taille.}
\begin{enumerate}
\item \question{Tester $H_0:\mu=5$ contre $H_1:\mu\neq5$, au seuil de $1\%$. Un non-rejet suffit-il à établir la conformité ?}
\indication{Calculer d'abord la moyenne et la variance corrigée ; le test porte sur une moyenne, pas sur chaque épaisseur.}
\reponse{$\bar x=4{,}97$, variance descriptive $0{,}0081$, variance corrigée $s^2=0{,}008265306$ et $s\simeq0{,}0909137$.
La statistique $Z=\frac{\overline X-5}{\frac{S}{\sqrt{50}}}$ est approximativement normale sous $H_0$ et $z_{\mathrm{obs}}=-2{,}33333$.
Comme $|z_{\mathrm{obs}}|<\Phi^{-1}(0{,}995)\simeq2{,}57583$, on ne rejette pas au seuil de $1\%$.
La $p$-valeur approchée vaut $0{,}01963$. Le non-rejet ne prouve pas l'égalité ni une conformité à une tolérance, qui demanderait de définir une marge admissible.}
\item \question{À partir de quel seuil le test normal approché rejette-t-il ?}
\reponse{Avec la convention de rejet $p_{\mathrm{val}}\leq\alpha$, il rejette pour $\alpha\geq0{,}01963$, soit environ $1{,}963\%$. C'est une analyse de sensibilité ; le seuil d'une étude doit être fixé avant les données.
Si l'on ajoute l'hypothèse de normalité des observations et utilise Student à 49 ddl, la $p$-valeur est environ $0{,}02378$ ; la décision au seuil de $1\%$ est identique.}
\end{enumerate}
}
